\section{Methods}

\subsection{Cohorts, harmonization and the no-imputation data layer}
We analysed the harmonized baseline assessment of $N=9{,}013$ patients with bipolar disorder
($6{,}252$), schizophrenia ($2{,}209$) or major depression ($552$) from three national FACE
cohorts across 21 expert centres. Instruments common to the cohorts were mapped to a shared
143-variable pool spanning cognition, symptom and functioning scales, sleep and chronotype,
developmental and substance-use history, and a routine biology panel (body composition, lipids,
glycaemia, urate, C-reactive protein). Missingness is pervasive and non-random in such registries;
rather than impute, we treat each patient's observed cells as the likelihood target and marginalize
the rest, so no cell is ever filled in. Diagnosis was withheld from model fitting throughout and
used only to describe, never to define, the map.

\subsection{The measurement model: a Bayesian bifactor structure}
The core is a single global Bayesian sparse bifactor / exploratory structural-equation
model~\citep{reise2012bifactor,asparouhov2009esem}. The generative hypothesis is that a small set of
latent axes is the \emph{only} reason a patient's many indicators move together: each patient $i$
carries an eight-dimensional latent score
\begin{equation}
f_i=(G_i,D_{i1},\dots,D_{iK})^{\top}\ \sim\ \mathcal{N}_{K+1}\!\left(\mathbf 0,\,\Phi\right),
\qquad \operatorname{diag}(\Phi)=\mathbf 1 ,
\label{eq:factors}
\end{equation}
a general functional-burden factor \Gfac{} and $K=7$ specific axes (cognition, immunometabolic load,
sleep, mania/activation, suicidality, developmental risk, substance use). $\Phi$ is a
\emph{correlation} matrix: its unit diagonal fixes the otherwise-arbitrary latent scale, so a factor's
variance is normalized to one and only the loadings carry magnitude. Every indicator $j$ is a linear
reflection of these factors through a single linear predictor
\begin{equation}
\eta_{ij}=\alpha_j+\lambda_{jG}\,G_i+\sum_{k=1}^{K}\lambda_{jk}\,D_{ik}+\beta_j^{\top}c_i ,
\label{eq:linpred}
\end{equation}
with item intercept $\alpha_j$, a \emph{general} loading $\lambda_{jG}$ on \Gfac, \emph{specific}
loadings $\lambda_{jk}$ (row $j$ of the loading matrix $\Lambda=[\lambda_G\mid\Lambda_D]$), and
item-level covariates $c_i$ that shift an indicator's mean (age, sex, site) but are never themselves
factor indicators. The term \emph{bifactor} is exactly this structure: every indicator may load on
\Gfac{} \emph{and} on its home specific factor, so general burden and a domain-specific signal are
separated within each item rather than confounded. It is \emph{exploratory} rather than confirmatory
because an indicator is \emph{allowed} small cross-loadings on plausibly related factors, governed by
a regularized-horseshoe prior~\citep{carvalho2009horseshoe,piironen2017horseshoe} that shrinks the
cross-loading matrix toward simple structure while leaving anchored home-factor loadings free---so
sparsity is chosen by the data rather than imposed.

\paragraph{Variance decomposition.} Taking the Gaussian indicators as the illustrative case (the
mixed-likelihood generalization follows below) and stacking one patient's items, $x_i=\mu_i+\lambda_G
G_i+\Lambda_D D_i+\varepsilon_i$---with mean term $\mu_i=\alpha+Bc_i$ collecting the intercepts and
covariate shifts of Eq.~\ref{eq:linpred}, and residual $\varepsilon_i\sim\mathcal N(\mathbf 0,\Psi)$,
diagonal $\Psi=\operatorname{diag}(\sigma_1^2,\dots,\sigma_J^2)$---each indicator's variance splits into
three interpretable parts,
\begin{equation}
\operatorname{Var}(x_{ij})=\underbrace{\lambda_{jG}^{2}}_{\text{general}}
+\underbrace{\textstyle\sum_{k=1}^{K}\lambda_{jk}^{2}}_{\text{specific}}
+\underbrace{\sigma_j^{2}}_{\text{unique}} ,
\label{eq:vardecomp}
\end{equation}
the share attributable to overall burden, to the domain axes, and to item-specific noise. The
\emph{defining hypothesis} that makes this decomposition valid is \textbf{local independence}: given a
patient's latent scores the items are conditionally independent,
$p(x_i\mid f_i,\Theta)=\prod_{j\in O_i}p(x_{ij}\mid f_i,\Theta)$, i.e.\ the factors explain away all
shared covariance. In the linear-Gaussian case this is \emph{identical} to ``$\Psi$ is diagonal.''
This is also why missing cells are marginalized, never imputed: an invented value would inject
between-cell covariance that the model would mistake for a factor.

\paragraph{Observation model and identification.} Indicators are generated through mixed likelihoods
appropriate to each variable's type---Gaussian for continuous and log-transformed laboratory values,
ordinal-logit for scales, Bernoulli for binary, negative-binomial for counts---with a Gaussian-copula
block coupling the correlated continuous biology panel~\citep{hui2017gllvm}, so no variable is coerced
onto a foreign scale. Because $f_i$ is unobserved it is integrated out, and fitting targets each
patient's observed cells only,
\begin{equation}
\log p\!\left(X_{\mathrm{obs}}\mid\Theta\right)=\sum_{i=1}^{N}\ \sum_{j\in O_i}\
\log F_j\!\left(X_{ij}\mid\eta_{ij},\psi_j\right),
\label{eq:obslik}
\end{equation}
where $\Theta$ collects all model parameters, $O_i$ is the set of indicators observed for patient $i$,
$F_j$ is the per-indicator likelihood chosen by variable type (above) and $\psi_j$ its
dispersion/threshold parameters. Unobserved indicators are marginalized via the Woodbury identity for
the copula block---which avoids repeatedly inverting a large patient-specific covariance submatrix and
is the
no-imputation invariant on which every downstream analysis rests. The bifactor identification fixes
\Gfac{} orthogonal to the specific axes ($\Phi$ block-diagonal); we treat this as an identifying
coordinate choice defining the \emph{primary map}, not an empirical claim about independence.

\subsection{The correlated-\texorpdfstring{\Gfac}{G} arm and the primary estimand}
To measure---rather than assume---whether biology is separable from severity, we re-estimated the
model with the block-diagonal constraint released, allowing \Gfac{} to correlate freely with the
specific axes:
\begin{equation}
\Phi=\begin{pmatrix}1 & \phi_{Gd}^{\top}\\ \phi_{Gd} & \Phi_D\end{pmatrix},\qquad
\phi_{Gd}=\operatorname{Corr}(\Gfac,D)\in[-1,1]^{7}.
\label{eq:corrg}
\end{equation}
The vector $\phi_{Gd}$ (Eq.~\ref{eq:corrg}) is the \emph{primary estimand} for the biology--severity
question; the bifactor map remains the coordinate system for all other analyses. Robustness of the
immunometabolic separation was assessed by re-estimating $\phi_{Gd}$ after adjusting each indicator
for age, sex, education and recruitment site, and then additionally for antipsychotic exposure.

\subsection{Estimation, invariance and scoring}
Posteriors were sampled with the No-U-Turn Sampler~\citep{hoffman2014nuts} in a probabilistic-%
programming implementation~\citep{abrilpla2023pymc,phan2019numpyro}, with a fast stochastic-%
variational re-estimation~\citep{blei2017vi,hoffman2013svi,kingma2014vae} used to confirm the atlas
at reduced cost (variational and sampled factor scores agree to within posterior uncertainty). Key
priors are weakly-informative: a regularized-horseshoe on the cross-loadings (global scale
$\tau_0=0.05$, slab $c=0.30$) and $\Phi\sim\mathrm{LKJ}(\eta{=}2)$. Intermediate stages advanced only
under a strict gate ($\hat R\le1.01$, effective sample size $\ge400$, zero divergences); the reported
joint eight-factor (mixed-likelihood) map is taken at the largest sample that mixes
($\hat R\le1.03$, zero divergences), with a cross-seed stability guard (Tucker congruence $\varphi=0.99$).
Measurement invariance was checked across the three cohorts (\edfig{edfig:invariance}) and by
leave-one-cohort-out refitting (\edfig{edfig:loso}), with split-half reproducibility benchmarking
(\edfig{edfig:repbench}); patient factor scores and their posterior uncertainty were obtained under
the frozen map. All validation analyses below were pre-specified against the fixed map;
the map was never re-fit to improve a downstream result.

\subsection{Structure testing and archetypes}
Whether the patient space is clustered or continuous was tested by comparing the best silhouette
(a standard cluster-separation score: within- versus between-cluster distance) over $K=2$--$12$
clusters against the distribution expected under a single-multivariate-Gaussian null (a parametric
bootstrap of the same covariance), and by Gaussian-mixture model selection. The continuum was
summarized by archetypal analysis~\citep{cutler1994archetypal}, which represents each patient as a
convex blend of extreme phenotypes (weights $w_a\ge 0$, $\sum_a w_a=1$); across candidate counts the
five-archetype solution was retained as the point beyond which reconstruction error flattened and
cross-seed archetype recovery remained stable. Descriptive tightness was quantified as the mean $\eta^{2}$ (the proportion of
between-patient variance in a factor score explained by an encoding---unrelated to the linear
predictor $\eta_{ij}$ above) of each encoding (archetype, DSM-5, cohort) across the eight axes: a
higher value means the encoding localizes a patient more tightly on the map.

\subsection{Temporal coherence and trait/state}
Two yearly follow-ups were scored under the frozen map. A between/within/measurement variance
decomposition yielded a test--retest intraclass correlation per axis,
\begin{equation}
\mathrm{ICC}=\frac{\tau^{2}}{\tau^{2}+\omega^{2}},
\label{eq:m-icc}
\end{equation}
with $\tau^{2}$ the durable between-patient (trait) variance and $\omega^{2}$ the within-patient
plus measurement (state) variance, separating durable traits from moving states.

\subsection{Prognosis: errors-in-variables incremental validity}
Two-year functional remission (based on the EGF, the French Global Assessment of Functioning scale)
and continuous functioning were modelled with an
errors-in-variables Bayesian generalized linear model that propagates each patient's factor-score
uncertainty into the outcome regression,
\begin{equation}
\theta_{id}\sim\mathcal N\!\big(\hat f_{id},\,s_{id}^{2}\big),\qquad
g\big(\mathbb E[y_i]\big)=\beta_0+\beta_W^{\top}W_i+\beta_\theta^{\top}\theta_i ,
\label{eq:m-eiv}
\end{equation}
with $\hat f_{id}$ the posterior factor score, $s_{id}$ its posterior SD and $W_i$ the covariates
(diagnosis, baseline severity, baseline outcome), avoiding regression dilution from treating latent
scores as known. Incremental validity was assessed by expected log pointwise predictive density on
held-out data~\citep{vehtari2017loo,watanabe2010waic}, with individual-level
discrimination (AUC) reported alongside. Because retention at two years is associated with baseline
map position, the incremental-validity estimate was confirmed under inverse-probability-of-attrition
weighting. The archetype-prognosis map (Fig.~\ref{fig:money}) bins
patients by their archetype-weight position and shades by mean two-year functional remission
($n=2{,}420$ with follow-up, bins $\ge 5$ patients).

\subsection{Data and code availability}
Individual patient data are governed by the FACE cohort agreements and are not publicly
redistributable; derived aggregate outputs underlying the figures, and the modelling and
figure-generation code, are available from the corresponding author on reasonable request. Per-%
patient data are never required to regenerate the aggregate results.
